%%%%%%%%%%%%%%%%%%%%%%%%
%%% 07-judges/11.tex %%%
%%%%%%%%%%%%%%%%%%%%%%%%

\Chapter{11}

\Verse{1}
{
	\hJ{וְיִפְתָּח הַגִּלְעָדִי הָיָה גִּבּוֹר חַיִל וְהוּא בֶּן־אִשָּׁה זוֹנָה וַיּוֹלֶד גִּלְעָד אֶת־יִפְתָּח׃}
}
{
	\eJ{
		Jephthah the Gileadite was an able warrior, who was the son of a
		prostitute. Jephthah’s father was Gilead;
	}
}
{
}

\Verse{2}
{
	\hJ{וַתֵּלֶד אֵשֶׁת־גִּלְעָד לוֹ בָּנִים וַיִּגְדְּלוּ בְנֵי־הָאִשָּׁה וַיְגָרְשׁוּ אֶת־יִפְתָּח וַיֹּאמְרוּ לוֹ לֹא־תִנְחַל בְּבֵית־אָבִינוּ כִּי בֶּן־אִשָּׁה אַחֶרֶת אָתָּה׃}
}
{
	\eJ{
		but Gilead also had sons by his wife, and when the wife’s sons
		grew up, they drove Jephthah out. They said to him, “You shall
		have no share in our father’s property, for you are the son of an
		outsider.”
	}
}
{
}

\Verse{3}
{
	\hJ{וַיִּבְרַח יִפְתָּח מִפְּנֵי אֶחָיו וַיֵּשֶׁב בְּאֶרֶץ טוֹב וַיִּתְלַקְּטוּ אֶל־יִפְתָּח אֲנָשִׁים רֵיקִים וַיֵּצְאוּ עִמּוֹ׃}
}
{
	\eJ{
		So Jephthah fled from his brothers and settled in the Tob country.
		Men of low character gathered about Jephthah and went out raiding
		with him.
	}
}
{
}

\Verse{4}
{
	\hJ{וַיְהִי מִיָּמִים וַיִּלָּחֲמוּ בְנֵי־עַמּוֹן עִם־יִשְׂרָאֵל׃}
}
{
	\eJ{Some time later, the Ammonites went to war against Israel.}
}
{
}

\Verse{5}
{
	\hJ{וַיְהִי כַּאֲשֶׁר־נִלְחֲמוּ בְנֵי־עַמּוֹן עִם־יִשְׂרָאֵל וַיֵּלְכוּ זִקְנֵי גִלְעָד לָקַחַת אֶת־יִפְתָּח מֵאֶרֶץ טוֹב׃}
}
{
	\eJ{
		And when the Ammonites attacked Israel, the elders of Gilead went
		to bring Jephthah back from the Tob country.
	}
}
{
}

\Verse{6}
{
	\hJ{וַיֹּאמְרוּ לְיִפְתָּח לְכָה וְהָיִיתָה לָּנוּ לְקָצִין וְנִלָּחֲמָה בִּבְנֵי עַמּוֹן׃}
}
{
	\eJ{
		They said to Jephthah, “Come be our chief, so that we can fight
		the Ammonites.”
	}
}
{
}

\Verse{7}
{
	\hJ{וַיֹּאמֶר יִפְתָּח לְזִקְנֵי גִלְעָד הֲלֹא אַתֶּם שְׂנֵאתֶם אוֹתִי וַתְּגָרְשׁוּנִי מִבֵּית אָבִי וּמַדּוּעַ בָּאתֶם אֵלַי עַתָּה כַּאֲשֶׁר צַר לָכֶם׃}
}
{
	\eJ{
		Jephthah replied to the elders of Gilead, “You are the very people
		who rejected me and drove me out of my father’s house. How can you
		come to me now when you are in trouble?”
	}
}
{
}

\Verse{8}
{
	\hJ{וַיֹּאמְרוּ זִקְנֵי גִלְעָד אֶל־יִפְתָּח לָכֵן עַתָּה שַׁבְנוּ אֵלֶיךָ וְהָלַכְתָּ עִמָּנוּ וְנִלְחַמְתָּ בִּבְנֵי עַמּוֹן וְהָיִיתָ לָּנוּ לְרֹאשׁ לְכֹל יֹשְׁבֵי גִלְעָד׃}
}
{
	\eJ{
		The elders of Gilead said to Jephthah, “Honestly, we have now
		turned back to you. If you come with us and fight the Ammonites,
		you shall be our commander over all the inhabitants of Gilead.”
	}
}
{
}

\Verse{9}
{
	\hJ{וַיֹּאמֶר יִפְתָּח אֶל־זִקְנֵי גִלְעָד אִם־מְשִׁיבִים אַתֶּם אוֹתִי לְהִלָּחֵם בִּבְנֵי עַמּוֹן וְנָתַן יְהֹוָה אוֹתָם לְפָנָי אָנֹכִי אֶהְיֶה לָכֶם לְרֹאשׁ׃}
}
{
	\eJ{
		Jephthah said to the elders of Gilead, “[Very well,] if you bring
		me back to fight the Ammonites and YHWH delivers them to me, I am
		to be your commander.”
	}
}
{
}

\Verse{10}
{
	\hJ{וַיֹּאמְרוּ זִקְנֵי־גִלְעָד אֶל־יִפְתָּח יְהֹוָה יִהְיֶה שֹׁמֵעַ בֵּינוֹתֵינוּ אִם־לֹא כִדְבָרְךָ כֵּן נַעֲשֶׂה׃}
}
{
	\eJ{
		And the elders of Gilead answered Jephthah, “YHWH Himself shall be
		witness between us: we will do just as you have said.”
	}
}
{
}

\Verse{11}
{
	\hJ{וַיֵּלֶךְ יִפְתָּח עִם־זִקְנֵי גִלְעָד וַיָּשִׂימוּ הָעָם אוֹתוֹ עֲלֵיהֶם לְרֹאשׁ וּלְקָצִין וַיְדַבֵּר יִפְתָּח אֶת־כּל־דְּבָרָיו לִפְנֵי יְהֹוָה בַּמִּצְפָּה׃}
}
{
	\eJ{
		Jephthah went with the elders of Gilead, and the people made him
		their commander and chief. And Jephthah repeated all these terms
		before YHWH at Mizpah.
	}
}
{
}

\Verse{12}
{
	\hJ{וַיִּשְׁלַח יִפְתָּח מַלְאָכִים אֶל־מֶלֶךְ בְּנֵי־עַמּוֹן לֵאמֹר מַה־לִּי וָלָךְ כִּי־בָאתָ אֵלַי לְהִלָּחֵם בְּאַרְצִי׃}
}
{
	\eJ{
		Jephthah then sent messengers to the king of the Ammonites,
		saying, “What have you against me that you have come to make war
		on my country?”
	}
}
{
}

\Verse{13}
{
	\hJ{
		וַיֹּאמֶר מֶלֶךְ בְּנֵי־עַמּוֹן אֶל־מַלְאֲכֵי יִפְתָּח כִּי־לָקַח יִשְׂרָאֵל אֶת־אַרְצִי בַּעֲלוֹתוֹ מִמִּצְרַיִם מֵאַרְנוֹן וְעַד־הַיַּבֹּק וְעַד־הַיַּרְדֵּן וְעַתָּה הָשִׁיבָה
		אֶתְהֶן בְּשָׁלוֹם׃
	}
}
{
	\eJ{
		The king of the Ammonites replied to Jephthah’s messengers, “When
		Israel came from Egypt, they seized the land which is mine, from
		the Arnon to the Jabbok as far as the Jordan. Now, then, restore
		it peaceably.”
	}
}
{
}

\Verse{14}
{
	\hJ{וַיּוֹסֶף עוֹד יִפְתָּח וַיִּשְׁלַח מַלְאָכִים אֶל־מֶלֶךְ בְּנֵי עַמּוֹן׃}
}
{
	\eJ{
		Jephthah again sent messengers to the king of the Ammonites.
	}
}
{
}

\Verse{15}
{
	\hJ{וַיֹּאמֶר לוֹ כֹּה אָמַר יִפְתָּח לֹא־לָקַח יִשְׂרָאֵל אֶת־אֶרֶץ מוֹאָב וְאֶת־אֶרֶץ בְּנֵי עַמּוֹן׃}
}
{
	\eJ{
		He said to him, “Thus said Jephthah: Israel did not seize the land
		of Moab or the land of the Ammonites.
	}
}
{
}

\Verse{16}
{
	\hJ{כִּי בַּעֲלוֹתָם מִמִּצְרָיִם וַיֵּלֶךְ יִשְׂרָאֵל בַּמִּדְבָּר עַד־יַם־סוּף וַיָּבֹא קָדֵשָׁה׃}
}
{
	\eJ{
		When they left Egypt, Israel traveled through the wilderness to
		the Sea of Reeds and went on to Kadesh.
	}
}
{
}

\Verse{17}
{
	\hJ{
		וַיִּשְׁלַח יִשְׂרָאֵל מַלְאָכִים אֶל־מֶלֶךְ אֱדוֹם לֵאמֹר אֶעְבְּרָה־נָּא בְאַרְצֶךָ וְלֹא שָׁמַע מֶלֶךְ אֱדוֹם וְגַם אֶל־מֶלֶךְ מוֹאָב שָׁלַח וְלֹא אָבָה וַיֵּשֶׁב יִשְׂרָאֵל
		בְּקָדֵשׁ׃
	}
}
{
	\eJ{
		Israel then sent messengers to the king of Edom, saying, ‘Allow us
		to cross your country.’ But the king of Edom would not consent.
		They also sent a mission to the king of Moab, and he refused. So
		Israel, after staying at Kadesh,
	}
}
{
}

\Verse{18}
{
	\hJ{
		וַיֵּלֶךְ בַּמִּדְבָּר וַיָּסב אֶת־אֶרֶץ אֱדוֹם וְאֶת־אֶרֶץ מוֹאָב וַיָּבֹא מִמִּזְרַח־שֶׁמֶשׁ לְאֶרֶץ מוֹאָב וַיַּחֲנוּן בְּעֵבֶר אַרְנוֹן וְלֹא־בָאוּ בִּגְבוּל מוֹאָב כִּי
		אַרְנוֹן גְּבוּל מוֹאָב׃
	}
}
{
	\eJ{
		traveled on through the wilderness, skirting the land of Edom and
		the land of Moab. They kept to the east of the land of Moab until
		they encamped on the other side of the Arnon; and, since Moab ends
		at the Arnon, they never entered Moabite territory.
	}
}
{
}

\Verse{19}
{
	\hJ{וַיִּשְׁלַח יִשְׂרָאֵל מַלְאָכִים אֶל־סִיחוֹן מֶלֶךְ־הָאֱמֹרִי מֶלֶךְ חֶשְׁבּוֹן וַיֹּאמֶר לוֹ יִשְׂרָאֵל נַעְבְּרָה־נָּא בְאַרְצְךָ עַד־מְקוֹמִי׃}
}
{
	\eJ{
		“Then Israel sent messengers to Sihon king of the Amorites, the
		king of Heshbon. Israel said to him, ‘Allow us to cross through
		your country to our homeland.’
	}
}
{
}

\Verse{20}
{
	\hJ{וְלֹא־הֶאֱמִין סִיחוֹן אֶת־יִשְׂרָאֵל עֲבֹר בִּגְבֻלוֹ וַיֶּאֱסֹף סִיחוֹן אֶת־כּל־עַמּוֹ וַיַּחֲנוּ בְּיָהְצָה וַיִּלָּחֶם עִם־יִשְׂרָאֵל׃}
}
{
	\eJ{
		But Sihon would not trust Israel to pass through his territory.
		Sihon mustered all his troops, and they encamped at Jahaz; he
		engaged Israel in battle.
	}
}
{
}

\Verse{21}
{
	\hJ{וַיִּתֵּן יְהֹוָה אֱלֹהֵי־יִשְׂרָאֵל אֶת־סִיחוֹן וְאֶת־כּל־עַמּוֹ בְּיַד יִשְׂרָאֵל וַיַּכּוּם וַיִּירַשׁ יִשְׂרָאֵל אֵת כּל־אֶרֶץ הָאֱמֹרִי יוֹשֵׁב הָאָרֶץ הַהִיא׃}
}
{
	\eJ{
		But YHWH, the God of Israel, delivered Sihon and all his troops
		into Israel’s hands, and they defeated them; and Israel took
		possession of all the land of the Amorites, the inhabitants of
		that land.
	}
}
{
}

\Verse{22}
{
	\hJ{וַיִּירְשׁוּ אֵת כּל־גְּבוּל הָאֱמֹרִי מֵאַרְנוֹן וְעַד־הַיַּבֹּק וּמִן־הַמִּדְבָּר וְעַד־הַיַּרְדֵּן׃}
}
{
	\eJ{
		Thus they possessed all the territory of the Amorites from the
		Arnon to the Jabbok and from the wilderness to the Jordan.
	}
}
{
}

\Verse{23}
{
	\hJ{וְעַתָּה יְהֹוָה אֱלֹהֵי יִשְׂרָאֵל הוֹרִישׁ אֶת־הָאֱמֹרִי מִפְּנֵי עַמּוֹ יִשְׂרָאֵל וְאַתָּה תִּירָשֶׁנּוּ׃}
}
{
	\eJ{
		“Now, then, YHWH, the God of Israel, dispossessed the Amorites
		before His people Israel; and should you possess their land?
	}
}
{
}

\Verse{24}
{
	\hJ{הֲלֹא אֵת אֲשֶׁר יוֹרִישְׁךָ כְּמוֹשׁ אֱלֹהֶיךָ אוֹתוֹ תִירָשׁ וְאֵת כּל־אֲשֶׁר הוֹרִישׁ יְהֹוָה אֱלֹהֵינוּ מִפָּנֵינוּ אוֹתוֹ נִירָשׁ׃}
}
{
	\eJ{
		Do you not hold what Chemosh your god gives you to possess? So we
		will hold on to everything that YHWH our God has given us to
		possess.
	}
}
{
}

\Verse{25}
{
	\hJ{וְעַתָּה הֲטוֹב טוֹב אַתָּה מִבָּלָק בֶּן־צִפּוֹר מֶלֶךְ מוֹאָב הֲרוֹב רָב עִם־יִשְׂרָאֵל אִם־נִלְחֹם נִלְחַם בָּם׃}
}
{
	\eJ{
		“Besides, are you any better than Balak son of Zippor, king of
		Moab? Did he start a quarrel with Israel or go to war with them?
	}
}
{
}

\Verse{26}
{
	\hJ{
		בְּשֶׁבֶת יִשְׂרָאֵל בְּחֶשְׁבּוֹן וּבִבְנוֹתֶיהָ וּבְעַרְעוֹר וּבִבְנוֹתֶיהָ וּבְכל־הֶעָרִים אֲשֶׁר עַל־יְדֵי אַרְנוֹן שְׁלֹשׁ מֵאוֹת שָׁנָה וּמַדּוּעַ לֹא־הִצַּלְתֶּם בָּעֵת
		הַהִיא׃
	}
}
{
	\eJ{
		“While Israel has been inhabiting Heshbon and its dependencies,
		and Aroer and its dependencies, and all the towns along the Arnon
		for three hundred years, why have you not tried to recover them
		all this time?
	}
}
{
}

\Verse{27}
{
	\hJ{וְאָנֹכִי לֹא־חָטָאתִי לָךְ וְאַתָּה עֹשֶׂה אִתִּי רָעָה לְהִלָּחֶם בִּי יִשְׁפֹּט יְהֹוָה הַשֹּׁפֵט הַיּוֹם בֵּין בְּנֵי יִשְׂרָאֵל וּבֵין בְּנֵי עַמּוֹן׃}
}
{
	\eJ{
		I have done you no wrong; yet you are doing me harm and making war
		on me. May YHWH, who judges, decide today between the Israelites
		and the Ammonites!”
	}
}
{
}

\Verse{28}
{
	\hJ{וְלֹא שָׁמַע מֶלֶךְ בְּנֵי עַמּוֹן אֶל־דִּבְרֵי יִפְתָּח אֲשֶׁר שָׁלַח אֵלָיו׃}
}
{
	\eJ{
		But the king of the Ammonites paid no heed to the message that
		Jephthah sent him.
	}
}
{
}

\Verse{29}
{
	\hJ{וַתְּהִי עַל־יִפְתָּח רוּחַ יְהֹוָה וַיַּעֲבֹר אֶת־הַגִּלְעָד וְאֶת־מְנַשֶּׁה וַיַּעֲבֹר אֶת־מִצְפֵּה גִלְעָד וּמִמִּצְפֵּה גִלְעָד עָבַר בְּנֵי עַמּוֹן׃}
}
{
	\eJ{
		Then the spirit of YHWH came upon Jephthah. He marched through
		Gilead and Manasseh, passing Mizpeh of Gilead; and from Mizpeh of
		Gilead he crossed over [to] the Ammonites.
	}
}
{
}

\Verse{30}
{
	\hJ{וַיִּדַּר יִפְתָּח נֶדֶר לַיהֹוָה וַיֹּאמַר אִם־נָתוֹן תִּתֵּן אֶת־בְּנֵי עַמּוֹן בְּיָדִי׃}
}
{
	\eJ{
		And Jephthah made the following vow to YHWH: “If you deliver the
		Ammonites into my hands,
	}
}
{
}

\Verse{31}
{
	\hJ{וְהָיָה הַיּוֹצֵא אֲשֶׁר יֵצֵא מִדַּלְתֵי בֵיתִי לִקְרָאתִי בְּשׁוּבִי בְשָׁלוֹם מִבְּנֵי עַמּוֹן וְהָיָה לַיהֹוָה וְהַעֲלִיתִיהוּ עֹלָה׃}
}
{
	\eJ{
		then whatever comes out of the door of my house to meet me on my
		safe return from the Ammonites shall be YHWH’s and shall be
		offered by me as a burnt offering.”
	}
}
{
%	\footnote{
%		Footnote 11:31. Jephthah's vow is worded with terrible openness: whoever comes
%		out. The tragedy is already inside the grammar.
%	}
}

\Verse{32}
{
	\hJ{וַיַּעֲבֹר יִפְתָּח אֶל־בְּנֵי עַמּוֹן לְהִלָּחֶם בָּם וַיִּתְּנֵם יְהֹוָה בְּיָדוֹ׃}
}
{
	\eJ{
		Jephthah crossed over to the Ammonites and attacked them, and YHWH
		delivered them into his hands.
	}
}
{
}

\Verse{33}
{
	\hJ{וַיַּכֵּם מֵעֲרוֹעֵר וְעַד־בֹּאֲךָ מִנִּית עֶשְׂרִים עִיר וְעַד אָבֵל כְּרָמִים מַכָּה גְּדוֹלָה מְאֹד וַיִּכָּנְעוּ בְּנֵי עַמּוֹן מִפְּנֵי בְּנֵי יִשְׂרָאֵל׃}
}
{
	\eJ{
		He utterly routed them—from Aroer as far as Minnith, twenty
		towns—all the way to Abel-cheramim. So the Ammonites submitted to
		the Israelites.
	}
}
{
}

\Verse{34}
{
	\hJ{וַיָּבֹא יִפְתָּח הַמִּצְפָּה אֶל־בֵּיתוֹ וְהִנֵּה בִתּוֹ יֹצֵאת לִקְרָאתוֹ בְּתֻפִּים וּבִמְחֹלוֹת וְרַק הִיא יְחִידָה אֵין־לוֹ מִמֶּנּוּ בֵּן אוֹ־בַת׃}
}
{
	\eJ{
		When Jephthah arrived at his home in Mizpah, there was his
		daughter coming out to meet him, with timbrel and dance! She was
		an only child; he had no other son or daughter.
	}
}
{
}

\Verse{35}
{
	\hJ{
		וַיְהִי כִרְאוֹתוֹ אוֹתָהּ וַיִּקְרַע אֶת־בְּגָדָיו וַיֹּאמֶר אֲהָהּ בִּתִּי הַכְרֵעַ הִכְרַעְתִּנִי וְאַתְּ הָיִית בְּעֹכְרָי וְאָנֹכִי פָּצִיתִי פִי אֶל־יְהֹוָה וְלֹא אוּכַל
		לָשׁוּב׃
	}
}
{
	\eJ{
		On seeing her, he rent his clothes and said, “Alas, daughter! You
		have brought me low; you have become my troubler! For I have
		uttered a vow to YHWH and I cannot retract.”
	}
}
{
}

\Verse{36}
{
	\hJ{וַתֹּאמֶר אֵלָיו אָבִי פָּצִיתָה אֶת־פִּיךָ אֶל־יְהֹוָה עֲשֵׂה לִי כַּאֲשֶׁר יָצָא מִפִּיךָ אַחֲרֵי אֲשֶׁר עָשָׂה לְךָ יְהֹוָה נְקָמוֹת מֵאֹיְבֶיךָ מִבְּנֵי עַמּוֹן׃}
}
{
	\eJ{
		“Father,” she said, “you have uttered a vow to YHWH; do to me as
		you have vowed, seeing that YHWH has vindicated you against your
		enemies, the Ammonites.”
	}
}
{
}

\Verse{37}
{
	\hJ{
		וַתֹּאמֶר אֶל־אָבִיהָ יֵעָשֶׂה לִּי הַדָּבָר הַזֶּה הַרְפֵּה מִמֶּנִּי שְׁנַיִם חֳדָשִׁים וְאֵלְכָה וְיָרַדְתִּי עַל־הֶהָרִים וְאֶבְכֶּה עַל־בְּתוּלַי אָנֹכִי (ורעיתי)
		[וְרֵעוֹתָי]׃
	}
}
{
	\eJ{
		She further said to her father, “Let this be done for me: let me
		be for two months, and I will go with my companions and lament
		upon the hills and there bewail my maidenhood.”
	}
}
{
	Qere Ketiv.
}

\Verse{38}
{
	\hJ{וַיֹּאמֶר לֵכִי וַיִּשְׁלַח אוֹתָהּ שְׁנֵי חֳדָשִׁים וַתֵּלֶךְ הִיא וְרֵעוֹתֶיהָ וַתֵּבְךְּ עַל־בְּתוּלֶיהָ עַל־הֶהָרִים׃}
}
{
	\eJ{
		“Go,” he replied. He let her go for two months, and she and her
		companions went and bewailed her maidenhood upon the hills.
	}
}
{
}

\Verse{39}
{
	\hJ{וַיְהִי מִקֵּץ שְׁנַיִם חֳדָשִׁים וַתָּשׁב אֶל־אָבִיהָ וַיַּעַשׂ לָהּ אֶת־נִדְרוֹ אֲשֶׁר נָדָר וְהִיא לֹא־יָדְעָה אִישׁ וַתְּהִי־חֹק בְּיִשְׂרָאֵל׃}
}
{
	\eJ{
		After two months’ time, she returned to her father, and he did to
		her as he had vowed. She had never known a man. So it became a
		custom in Israel
	}
}
{
}

\Verse{40}
{
	\hJ{מִיָּמִים יָמִימָה תֵּלַכְנָה בְּנוֹת יִשְׂרָאֵל לְתַנּוֹת לְבַת־יִפְתָּח הַגִּלְעָדִי אַרְבַּעַת יָמִים בַּשָּׁנָה׃}
}
{
	\eJ{
		for the maidens of Israel to go every year, for four days in the
		year, and chant dirges for the daughter of Jephthah the Gileadite.
	}
}
{
}
